\scanpage{52}

Étale algebras and the different (4)\hfill Tribu number 65

§ 2. \uline{Different.}

\uline{number 1.}\quad \uline{Filting ideals of a module} $M$.

See D. Lazard's thesis. Link with the content. The first Filting lies between $\operatorname{Ann}(M)$ and one of its powers.

\uline{number 2.}\quad \uline{Kähler different.}

We take for $B$ an algebra of \margin{\figure{\begin{tikzpicture}[baseline=(e.base)]
\node[anchor=base,inner sep=0pt] (e) at (0.50,0) {!};
\draw[line width=0.45pt] (0,0.02) -- (0.23,0.02);
\end{tikzpicture}}}finite presentation over $A$ (more generally perhaps it suffices to assume that $\Omega_A(B)$ is a $B$-module of finite presentation). By definition the Kähler different $\uline{D}_{B/A}$ is the first Filting ideal of $\Omega_A(B)$ :

if
\diagram{\begin{tikzcd}[ampersand replacement=\&,column sep=small]
A^p\arrow[r,"u"]\&A^q\arrow[r]\&\Omega_A(B)\arrow[r]\&0
\end{tikzcd}}
, it is $\operatorname{Im}(\Lambda^q u)$.

The zero locus of $\uline{D}_{B/A}$ is the set of ramification points. Behavior of $\uline{D}_{B/A}$ under change of ring (this is trivial).

\uline{number 3.}\quad \uline{Transitivity.}

Say what is known. If we have $A\to B\to C$ with $C$ a local complete intersection over $B$, the transitivity formula works.

\uline{number 4.}\quad \uline{Link with the trace.}

Define the Dedekind different by the trace, and show that it coincides with that of Kähler when $A$ and $B$ are discrete valuation rings (this is the big theorem 3 of number 7, § 2, with its dévissage, - possibly embellished with the lemmas of number 6 if they have not been done in an earlier chapter to illustrate the notion of a complete intersection. In any case we shall have said that if $A$ and $B$ are regular, $A\to B$ is a relative complete intersection. Proposition 11 will be clarified by the following lemma :

``If $k\to R$ is finite and if $R/\operatorname{rad}(R)$ is étale, then $R$ monogenic if and only if the maximal ideals of $R$ are homogeneous''.

Proposition 12 is absorbed by the above result on regular rings.

\uline{number 5.}\quad \uline{Complements.}

The tame ramification of number 5 (note that the reduction formula for traces used at the top of page 25 is done in Algebra VIII, § 11).

Give the Kummer description of totally and tamely ramified extensions (extract the n-th root of a well-chosen uniformizer ; this in the complete case); example of Riemann surfaces (here the Galois group has a canonical generator $\exp\left(\dfrac{2i\pi}{n}\right)$).

\mspage{51}
